\documentclass[../../main.tex]{subfiles}
\begin{document}

{\bf [解]}

題意より$x^{n-1}$次関数$P_n(x)$として
\begin{align}
  x^{n+1} = (x^2-x-1) P_n(x) + a_n x + b_n \label{eq:1}
\end{align}
とおく．

\paragraph{(1)}

\cref{eq:1}の両辺に $x$ をかけて
\begin{align}
  x^{n+2}
  &= x(x^2-x-1) P_n(x) + a_n x^2 + b_n x \\
  &= (x^2-x-1) [x P_n(x) + a_n] + (a_n + b_n) x + a_n \label{eq:2}
\end{align}
である．一方で\cref{eq:1}で$n\to n+1$として
\begin{align}
  x^{n+2} = (x^2-x-1) P_{n+1}(x) + a_{n+1} x + b_{n+1} \label{eq:3}
\end{align}
である．従って\cref{eq:2,eq:3}で係数を比較して
\begin{align}
  \begin{cases}
    a_{n+1} = a_n + b_n \\
    b_{n+1} = a_n
  \end{cases} \label{eq:4}
\end{align}
である．よって題意は示された．

\paragraph{(2)}

\cref{eq:1}で$n=1$として
\begin{align}
  x^2 = (x^2-x-1) + x + 1
\end{align}
だから，初項は
\begin{align}
  a_1 = b_1 = 1 \label{eq:5}
\end{align}
である．\cref{eq:4}から帰納的に$a_n, b_n \in \mathbb{Z}_{>0}$ である．

そこで以下$a_n$と$b_n$が互いに素であることを示す．
数学的帰納法によって題意を示す．

$n=1$の時，\cref{eq:5}より$(a_n,b_n)$は互いに素である．

次に$k\in\mathbb{N}$に対して$a_{k}$と$b_{k}$が互いに素であると仮定する．
\cref{eq:4}より「$a_{k+1}$ と $b_{k+1}$ が互いに素」$iff$ 「$b_{k}$ と $a_{k}$ が互いに素」だから，仮定より$a_{k+1}$ と $b_{k+1}$も互いに素である．

以上から数学的帰納法により任意の$n$について$(a_n,b_n)$は互いに素である．
よって題意は示された．

\end{document}
